\worldchapter{Der Charakter}{character}
\index{Charakter}
\label{ch:character}

\begin{multicols}{2}

Um die Welt von \trans{book_title} zu betreten, benötigst du einen Charakter – eine fiktive Person, der du Leben einhauchst und deren Schicksal du lenkst. Auf dem Charakterbogen werden all die Werte und Eigenschaften festgehalten, die deinen Helden einzigartig machen.

Dieses Kapitel gibt einen Überblick über die grundlegenden Eigenschaften eines Charakters. Der Prozess der Charaktererschaffung wird in \cref{ch:create} genau beschrieben.

\section{Persona}
\index{Persona}

Die Werte unter Persona beziehen sich auf die mentalen Eigenschaften des Charakters. Jeder Wert entspricht einer Persönlichkeitseigenschaft. Die Persona-Eigenschaften haben jeweils einen eigenen Wert und bilden auch die Grundlage für die Fertigkeiten.

\subsection{Bildung}
\index{Attribut!Bildung}\index{Bildung}

Bildung misst das angeeignete Allgemeinwissen und die Fähigkeit eines Charakters, auf erlernte Informationen zurückzugreifen. Ein hoher Wert deutet auf eine Person hin, die viel gelesen hat, eine gute Schule besuchte oder eine unstillbare Neugier besitzt. Folglich glänzt ein solcher Charakter in allen theoretischen Fertigkeiten wie \textit{Naturkunde} oder \textit{Historie}.

\subsection{Logik}
\index{Attribut!Logik}\index{Logik}

Während \textit{Bildung} das gespeicherte Wissen eines Charakters darstellt, ist Logik seine Fähigkeit, dieses Wissen anzuwenden und neue Schlüsse zu ziehen. Logik ist immer dann entscheidend, wenn es darum geht, aus vorhandenen Hinweisen ein stimmiges Gesamtbild zu erschaffen oder komplexe Probleme zu analysieren. Ein Charakter mit hoher Logik ist daher in Fertigkeiten wie \textit{Untersuchen} überlegen, um an einem Tatort die entscheidenden Details zu verknüpfen, oder in \textit{Mechanik}, um den Aufbau einer Falle zu durchschauen und sie zu entschärfen.

\subsection{Gewissenhaftigkeit}
\index{Attribut!Gewissenhaftigkeit}\index{Gewissenhaftigkeit}

Gewissenhaftigkeit beschreibt die Sorgfalt, Disziplin und Verlässlichkeit eines Charakters. Ein hoher Wert steht für methodisches Vorgehen und verhindert Flüchtigkeitsfehler, was bei Fertigkeiten wie \textit{Erste Hilfe} oder \textit{Heimlichkeit} entscheidend ist.

\subsection{Willenskraft}
\index{Attribut!Willenskraft}\index{Willenskraft}

Willenskraft ist die mentale Stärke und Entschlossenheit, mit der ein Charakter seine Ziele verfolgt und Widrigkeiten widersteht. Sie repräsentiert die innere Zähigkeit, sich weder von äußerem Druck noch von inneren Zweifeln von seinem Weg abbringen zu lassen. Ein hoher Wert ist daher die Grundlage für Fertigkeiten wie \textit{Tapferkeit}, um auch in aussichtslosen Lagen standhaft zu bleiben, oder \textit{Einschüchtern}, um den eigenen Willen auf andere zu projizieren.

\subsection{Auffassungsgabe}
\index{Attribut!Auffassungsgabe}\index{Auffassungsgabe}

Die Auffassungsgabe beschreibt, wie schnell und präzise ein Charakter seine Umgebung mit allen Sinnen erfasst und verarbeitet. Sie ist das Maß für die angeborene Aufmerksamkeit und das Bewusstsein für die Umgebung, vom kleinsten Geräusch bis zur weitesten Landschaft. Ein hoher Wert ist daher die Grundlage für Fertigkeiten wie \textit{Wahrnehmung}, um versteckte Details oder Gefahren zu entdecken, und \textit{Orientierung}, um sich in fremdem Terrain nicht zu verirren.

\subsection{Charme}
\index{Attribut!Charme}\index{Charme}

Charme ist die Fähigkeit, durch Wortwahl, Auftreten und Persönlichkeit eine positive Verbindung zu anderen herzustellen und sie für sich zu gewinnen. Im Gegensatz zur rein äußerlichen \textit{Attraktivität} ist Charme ein Attribut, das auf sozialer Intelligenz beruht. Ein hoher Wert in diesem Bereich ist die Basis für Fertigkeiten wie \textit{Politik}, um Verbündete zu gewinnen und Verhandlungen zu führen, sowie \textit{Empathie}, um die Stimmungen anderer zu verstehen und darauf zu reagieren.

\section{Physis}
\index{Physis}

Alle Physis-Attribute beschreiben die körperlichen Fähigkeiten des Charakters. Jedes Attribut hat einen Wert, der die Anzahl der Würfel angibt, die für dieses Attribut geworfen werden.

\subsection{Geschick}
\index{Attribut!Geschick}\index{Geschick}

Geschick umfasst sowohl die allgemeine Körperbeherrschung als auch die feine Hand-Augen-Koordination eines Charakters. Es ist das Maß für Reflexe, Balance und die Fähigkeit, präzise und kontrollierte Bewegungen auszuführen. Ein hoher Wert in diesem Attribut ermöglicht einem Charakter, in Fertigkeiten wie \textit{Akrobatik} zu brillieren, um Hindernissen auszuweichen und das Gleichgewicht zu halten, sowie im \textit{Schießen}, um ein Ziel auch auf große Entfernung sicher zu treffen.

\subsection{Kraft}
\index{Attribut!Kraft}\index{Kraft}

Kraft ist das Maß für die rohe Muskelkraft und die körperliche Stärke eines Charakters. Sie repräsentiert die Fähigkeit, überwältigende körperliche Gewalt auszuüben, sei es, um schwere Objekte zu bewegen oder um im Kampf verheerenden Schaden anzurichten. Ein hoher Kraftwert ist daher entscheidend für den \textit{Nahkampf}, um mit Hieben eine Rüstung zu durchschlagen, sowie für das \textit{Werfen}, um Objekte mit großer Wucht auf ein Ziel zu schleudern.

\subsection{Attraktivität}
\index{Attribut!Attraktivität}\index{Attraktivität}

Attraktivität misst die unmittelbare Wirkung der körperlichen Erscheinung und Präsenz eines Charakters auf andere. Dieser Wert beschreibt nicht zwangsläufig klassische Schönheit, sondern vielmehr, wie einprägsam oder fesselnd das Äußere einer Person ist - sei es durch anmutige Züge, eine einschüchternde Statur oder markante Merkmale. Eine hohe Attraktivität sorgt dafür, dass ein Charakter aus der Masse heraussticht und einen starken ersten Eindruck hinterlässt, noch bevor er ein Wort gesprochen hat.

\subsection{Ausdauer}
\index{Attribut!Ausdauer}\index{Ausdauer}

Ausdauer beschreibt die rein körperliche Widerstandsfähigkeit und Zähigkeit eines Charakters. Sie bestimmt, wie lange jemand anstrengenden Tätigkeiten wie einem Gewaltmarsch, einem Dauerlauf oder einem Kampf nachgehen kann, bevor Erschöpfung einsetzt.

\subsection{Resistenz}
\index{Attribut!Resistenz}\index{Resistenz}

Resistenz ist die angeborene Zähigkeit und Konstitution des Körpers, die es einem Charakter erlaubt, Schaden und feindlichen Einflüssen zu widerstehen. Sie repräsentiert die Fähigkeit, die Wirkung von Giften zu mindern oder den Strapazen extremer Hitze und Kälte zu trotzen. Ein Charakter mit hoher Resistenz kann Krankheiten und anderen schädlichen Effekten besser widerstehen.

\subsection{Schnelligkeit}
\index{Attribut!Schnelligkeit}\index{Schnelligkeit}

Schnelligkeit misst sowohl die reine Bewegungsgeschwindigkeit als auch die Reaktionsgabe eines Charakters. Sie bestimmt, wie rasant sich eine Person fortbewegen kann, um Distanzen zu überbrücken, aber auch, wie schnell sie auf plötzliche Ereignisse oder Gefahren reagieren kann.

\section{Der Mindestwurf}
\index{Mindestwurf}

Der Mindestwurf ist ein zentrales Merkmal des Charakters. Er gibt an, welches Ergebnis ein Würfel haben muss, um einen Erfolg darzustellen. Der Mindestwurf wird von der Abstammung übernommen und liegt bei den meisten Charakteren bei 5+. Die Schablone ``Meisterliche Präsenz'' senkt den Mindestwurf um 1, ansonsten kann er nur durch besondere Ereignisse oder seltene Gegenstände und oft nur für kurze Zeit verändert werden.

\section{Ausweichen}
\index{Ausweichen}

\textit{Ausweichen} wird im Kampf verwendet und erlaubt es, einem Nahkampfangriff auszuweichen. Der Wert ist gleich dem Ausweichwert der Abstammung plus dem Durchschnitt von \textit{Schnelligkeit} und \textit{Geschick} (aufgerundet). Rüstung und Waffen reduzieren diesen Wert. Charakterschablonen können diesen Wert verändern.

\section{Schutz}
\index{Schutz}

Hat ein Charakter aufgrund seiner Abstammung Schutz, so gilt dieser als ``angeborener Schutz''. Im Gegensatz zu den anderen Schutzarten bleibt dieser nicht bis zum Ende des Kampfes aufgebraucht, sondern wird zu Beginn jeder Kampfrunde des Spielers erneuert (siehe \cref{ch:combat}).

\section{Zusatzwürfel}
\index{Zusatzwürfel}

Jeder Charakter kann eine Anzahl von \textit{Bonuswürfeln}, \textit{Schicksalswürfeln} oder \textit{Wiederholungswürfeln} haben. Alle drei haben unterschiedliche Anwendungen (siehe \cref{ch:rolls}), stellen aber immer einen Vorteil für den Charakter dar, der während des Spiels genutzt werden kann.

Der Charakter kann verbrauchte Würfel während der Rast zurückerhalten (siehe \cref{ch:wounds}).

\section{Fertigkeiten}
\index{Fertigkeiten}

Komplexere Handlungen oder Kenntnisse werden durch \textit{Fertigkeiten} beschrieben. Alle Charaktere haben die gleichen Fertigkeiten mit unterschiedlichen Werten, so dass die Spielleiterin sicher sein kann, dass ein Spieler auf jeden Fall eine Fertigkeit würfeln kann.

Jede Fertigkeit besteht aus einem Basisattribut und einem Fertigkeitswert. Zum Beispiel ist der Basiswert für die Fertigkeit \textit{Einschüchtern} der Attributswert \textit{Auffassungsgabe}.

Zum Basisattributwert werden die Boni der gewählten Charakterschablonen addiert.

\subsection{Einschüchtern}
\index{Fertigkeit!Einschüchtern}\index{Einschüchtern}

Die Fertigkeit Einschüchtern stellt die Fähigkeit des Charakters dar, andere einzuschüchtern. Diese Fertigkeit kann z.B. eingesetzt werden, um Informationen aus einem Gegner herauszupressen oder Gegner dazu zu bringen, sich aus einem Kampf zurückzuziehen.

Attribut: \textit{Auffassungsgabe}

\subsection{Empathie}
\index{Fertigkeit!Empathie}\index{Empathie}

Empathie ist die Fähigkeit, Gefühle und Stimmungen eines Gegenübers zu deuten, und vielleicht Intentionen zu erkennen. Hierbei können keine Gedanken gelesen werden.

Attribut: \textit{Charme}

\subsection{Heimlichkeit}
\index{Fertigkeit!Heimlichkeit}\index{Heimlichkeit}

Heimlichkeit ist die Kunst des Verborgenen. Diese Fertigkeit wird sowohl für das heimliche Bewegen (Schleichen) verwendet, als auch um zu prüfen wie verschlossen der Charakter ist. Sie kann beispielsweise angewendet werden wenn der Charakter verhört wird oder versucht ist, ein Geheimnis auszuplaudern.

Attribut: \textit{Gewissenhaftigkeit}

\wbif{nexus}{}{
\subsection{Magiekunde}
\index{Fertigkeit!Magiekunde}\index{Magiekunde}

\wbif{phasesix}{Diese Fertigkeit ist nur in Kampagnen mit Magie verfügbar.}{}

Magiekunde ist das theoretische und akademische Wissen über die Natur der Magie. Sie ermöglicht es einem Charakter, die Struktur eines fremden Zaubers zu analysieren, arkane Schriften zu entziffern oder die Geschichte und Rituale alter magischer Kulte zu kennen. Im direkten Gegensatz zum praktischen \textit{Zaubern} ist Magiekunde die Fertigkeit des Verstehens und Identifizierens, nicht des Anwendens.
}

Attribut: \textit{Charme}

\subsection{Orientierung}
\index{Fertigkeit!Orientierung}\index{Orientierung}

Diese Fertigkeit dient der Orientierung, sowohl in der Landschaft, als auch in unübersichtlichen Situationen. So kann sie im unübersichtlichen Gedränge der Stadt angewandt werden, aber auch wenn der Charakter zum Beispiel durch einen Wasserstrudel durchgewirbelt wird.

Attribut: \textit{Auffassungsgabe}

\subsection{Politik}
\index{Fertigkeit!Politik}\index{Politik}

Immer wenn es darum geht, politisches Handeln abzuschätzen, wird diese Fertigkeit verwendet.
\wbif{tirakan}{Dies kann in der realen Politik der Fall sein, aber auch auf das sichere Bewegen in Adelskreisen darstellen.}{Dies kann in der realen Politik der Fall sein, aber auch auf das sichere Bewegen in großen Konzernen darstellen.}

Attribut: \textit{Charme}

\subsection{Religion}
\index{Fertigkeit!Religion}\index{Religion}

Diese Fertigkeit umfasst die Kenntnis von religiösen Lehren, aber auch die Sicherheit bei der Durchführung religiöser Zeremonien.

Attribut: \textit{Gewissenhaftigkeit}

\subsection{Tapferkeit}
\index{Fertigkeit!Tapferkeit}\index{Tapferkeit}

Diese Fertigkeit kommt immer dann zum Tragen, wenn es darum geht wie Mutig ein Charakter ist. Sie kann zum Beispiel verwendet werden wenn es darum geht, ob ein Charakter tapfer genug ist, sich einem übermächtigen Gegner gegenüberzustellen.

Attribut: \textit{Willenskraft}

\subsection{Täuschen und Betrügen}
\index{Fertigkeit!Täuschen und Betrügen}\index{Täuschen und Betrügen}

Wenn der Charakter ein Gegenüber täuschen, oder zum Beispiel beim Spiel betrügen möchte, kann auf diese Fertigkeit geworfen werden.

Attribut: \textit{Charme}

\subsection{Überzeugen}
\index{Fertigkeit!Überzeugen}\index{Überzeugen}

Will der Charakter sein Gegenüber argumentativ überzeugen, so kommt diese Fertigeit zur Anwendung.

Attribut: \textit{Willenskraft}

\subsection{Untersuchen}
\index{Fertigkeit!Untersuchen}\index{Untersuchen}

Diese Fertigkeit wird verwendet, wenn der Charakter einen Gegenstand, eine bestimmte Szene oder ein Objekt auf bestimmte Eigenschaften untersuchen möchte.

Attribut: \textit{Auffassungsgabe}

\subsection{Wahrnehmung}
\index{Fertigkeit!Wahrnehmung}\index{Wahrnehmung}

Wahrnehmung stellt die Fähigkeit des Charakters dar, Dinge in seinem Umfeld wahrzunehmen. Dies kann das Durchsuchen eines Hauses, die Suche nach dem schemenhaften Dieb am Waldrand, oder auch eine Regung im Gesicht des Gegenübers sein.

Attribut: \textit{Auffassungsgabe}

\subsection{Akrobatik}
\index{Fertigkeit!Akrobatik}\index{Akrobatik}

Akrobatik ist die Kunst, sich schnell und geschickt zu bewegen. Diese Fertigkeit kommt zum Einsatz, wenn der Charakter über einen Vorsprung klettert oder einen kurzen Sprint hinlegt. Der Wert kann auch für das Balancieren verwendet werden.

Attribut: \textit{Geschick}

\subsection{Darbietung}
\index{Fertigkeit!Darbietung}\index{Darbietung}

Darbietung ist das künstlerische Auftreten. Dies kann Schauspiel, aber auch der musikalische Vortrag eines Stücks sein. Eine imposante Lügengeschichte kann auch mit Hilfe der Darbietung erzählt werden.

Attribut: \textit{Charme}

\subsection{Erste Hilfe}
\index{Fertigkeit!Erste Hilfe}\index{Erste Hilfe}

Die Erste Hilfe muss mit ausreichend Verbandsmaterial durchgeführt werden, um erfolgreich zu sein.

Bei einem erfolgreichen Wurf werden die Wunden des Verletzten um die Hälfte (aufgerundet) der Erfolge des Wurfes geheilt. Durch die Erste Hilfe werden Blutungen gestillt.

Attribut: \textit{Gewissenhaftigkeit}

\subsection{Fahren}
\index{Fertigkeit!Fahren}\index{Fahren}

Die Fertigkeit Fahren beschreibt das Führen von Fahrzeugen aller Art. Hierbei trifft die Fertigkeit auf alle fahrbaren Objekte wie etwa Schiffe, Fahrzeuge oder Kutschen zu.

Attribut: \textit{Geschick}

\subsection{Historie}
\index{Fertigkeit!Historie}\index{Historie}

Historie beschreibt das Wissen des Charakters um Geschichte und vergangene Ereignisse. Auch Altertümer können mit dieser Fertigkeit eingeschätzt werden.

Attribut: \textit{Bildung}

\subsection{Kommunikation}
\index{Fertigkeit!Kommunikation}\index{Kommunikation}

Die Fähigkeit zur gesellschaftlichen Unterhaltung wird mit der Fertigkeit Kommunikation beschrieben. Sie beschreibt wie geschickt sich der Charakter in Unterhaltungen verhält.

Attribut: \textit{Bildung}

\subsection{Mechanik}
\index{Fertigkeit!Mechanik}\index{Mechanik}

Mechanik fasst alle handwerklichen Tätigkeiten, sowie die Kenntnis um mechanische Abläufe zusammen. Das Bearbeiten eines Holzstücks oder das Verstehen einer mechanischen Uhr können mit dieser Fertigkeit abgebildet werden.

Attribut: \textit{Logik}

\subsection{Nahkampf}
\index{Fertigkeit!Nahkampf}\index{Nahkampf}

Der Wert dieser Fertigkeit stellt die Grundlage für den Angriff mit Nahkampfwaffen dar. Auf diese Fertigkeit wird in der Regel nicht direkt geworfen.

Attribut: \textit{Kraft}

\subsection{Naturkunde}
\index{Fertigkeit!Naturkunde}\index{Naturkunde}

Naturkunde beschreibt die Kenntnis des Charakters von allen Facetten der Natur. Diese Fertigkeit kann zur Anwendung kommen wenn der Charakter nach Pflanzen sucht, im Wald Holz sammelt oder das Wesen eines Tiers einschätzen will.

Attribut: \textit{Bildung}

\subsection{Schießen}
\index{Fertigkeit!Schießen}\index{Schießen}

Der Wert dieser Fertigkeit stellt die Grundlage für den Angriff mit Fernkampfwaffen dar. Auf diese Fertigkeit wird in der Regel nicht direkt geworfen.

Attribut: \textit{Geschick}

\subsection{Werfen}
\index{Fertigkeit!Werfen}\index{Werfen}

Diese Fertigkeit wird immer dann verwendet, wenn der Charakter Gegenstände wirft. Dies können einfache Gegenstände wie Steine, aber auch Brandsätze oder Netze sein.

Für genaue Regeln für das Werfen von Gegenständen siehe \cref{sec:throwing}.

Attribut: \textit{Kraft}

\wbif{nexus}{}{
\subsection{Zaubern}
\index{Fertigkeit!Zaubern}\index{Zaubern}

\wbif{phasesix}{Diese Fertigkeit ist nur in Kampagnen mit Magie verfügbar.}{}

Die Fertigkeit Zaubern ist das Maß für die Beherrschung roher magischer Energien und die Fähigkeit, diese gezielt zu formen und zu lenken. Mit ihr werden mächtige Zauber gewirkt, magische Duelle ausgetragen oder Artefakte erschaffen. Im entscheidenden Gegensatz zur rein theoretischen \textit{Magiekunde} beschreibt Zaubern somit den aktiven und praktischen Einsatz von Magie.

Attribut: \textit{Willenskraft}
}

\section{Wissen}
\index{Wissen}

Wissen funktioniert ähnlich wie Fertigkeiten, allerdings ist hier die Liste nicht vorgegeben. Charaktere können aufgrund ihres Werdegangs verschiedene Wissensfertigkeiten haben, die sie frei einsetzen können. Wissen wird jeweils einer Fertigkeit zugeordnet. Der effektive Wert, auf den gewürfelt wird, setzt sich aus dem Wissenswert und dem Fertigkeitswert zusammen.

Wissen wird durch Charakterschablonen erlangt. Auf den Charakterschablonen ist angegeben, ob diese Wissen mitbringen.

\section{Schatten}
\index{Schatten}

Ein Charakter kann besondere Eigenschaften haben, die ihn außerhalb seiner körperlichen oder psychischen Eigenschaften beeinflussen. Jeder \textit{Schatten} hat seine eigene Beschreibung oder Regel. So kann ein Charakter z.B. einen Rivalen haben oder obrigkeitshörig sein. Schatten haben keine Werte, können aber ihre eigenen Regeln mitbringen.

Schatten sind auf Charakterschablonenen angegeben. Enthält eine Charakterschablone eine ausgeschriebene Regel, so ist dies ein Schatten.

\section{Sprachen}
\index{Sprachen}

Ein Charakter kann eine bestimmte Anzahl von Sprachen entsprechend der Summe seiner Attribute \textit{Bildung} und \textit{Logik} erlernen. Dies können beliebige Sprachen der entsprechenden Welt sein. Beträgt die Summe der Attribute 0 oder weniger, beherrscht der Charakter selbst seine Muttersprache nur eingeschränkt.

Der Grenzwert für die Anzahl der erlernten Sprachen dient als Richtwert für neue Charaktere. Im Spiel erlernte Sprachen können selbstverständlich auch über diesen Wert hinaus aufgeführt werden.

\wbif{nexus}{
Charakterschablonen oder Körpermodifikationen können die Anzahl der erlernbaren Sprachen erhöhen.
}{}
\wbif{tirakan}{
Charakterschablonen oder magische Gegenstände können die Anzahl der erlernbaren Sprachen erhöhen.
}{}
\wbif{phasesix}{
Charakterschablonen, Körpermodifikationen oder magische Gegenstände können die Anzahl der erlernbaren Sprachen erhöhen.
}{}

\section{Kontakte}
\index{Kontakte}

Kontakte sind Verbindungen des Charakters zu anderen Menschen oder Wesen, auf die er sich verlassen kann. In der Regel sind dies Kontakte außerhalb der Spielrunde, beispielsweise ein Adliger, ein Kontakt in der Regierung oder ein Arzt.

Bei der Charaktererschaffung kann der Charakter eine bestimmte Anzahl an Kontakten haben, die sich nach der Summe der Attribute \textit{Charme} und \textit{Attraktivität} richtet.

Natürlich kann diese Anzahl überschritten werden, wenn im Spiel neue Kontakte geschlossen werden.
\end{multicols}
